% from ega1-10-fr.tex
\oldpage[I]{208}

  \item[c)] Tout point de $\mathfrak{Y}$ possède un voisinage ouvert formel
    affine noethérien $V$ ayant la propriété suivante~:
    \begin{enumerate}
      \item[(\textbf{Q})] $f^{-1}(V)$ est réunion d'une famille finie
        d'ouverts formels affines noethériens $U_i$ tels que l'anneau adique
        noethérien $\Gamma(U_i,\mathscr{O}_{\mathfrak{X}})$ soit
        topologiquement isomorphe au quotient d'une algèbre de séries
        formelles restreintes (\hyperref[0.7.5.1-fr]{0, 7.5.1}) sur
        $\Gamma(V,\mathscr{O}_{\mathfrak{Y}})$, par un idéal (nécessairement
        fermé).
    \end{enumerate}
\end{enumerate}
\end{proposition}

Il est immédiat que a) entraîne b) en vertu de
(\hyperref[I.10.12.3-fr]{10.12.3}). Pour montrer que b) entraîne c), on
peut, puisque la question est locale sur $\mathfrak{Y}$, supposer que
$\mathfrak{Y}=\operatorname{Spf}(B)$, où $B$ est adique noethérien~; soit
$\mathscr{K}=\mathfrak{K}^\Delta$, $\mathfrak{K}$ étant un idéal de
définition de $B$. Comme par hypothèse, $X_0$ est de type fini sur $Y_0$,
$X_0$ est réunion finie d'ouverts affines $U_i$ tels que l'anneau $A_{i0}$
du schéma affine induit par $X_0$ sur $U_i$ soit une algèbre de type fini sur
l'anneau $B/\mathfrak{K}$ de $Y_0$
(\hyperref[I.6.3.2-fr]{6.3.2}). En vertu de
(\hyperref[I.5.1.9-fr]{5.1.9}), $U_i$ est aussi un ouvert affine dans chacun
des préschémas noethériens $X_n$, et si $A_{in}$ est l'anneau du schéma
affine induit par $X_n$ sur $U_i$, l'hypothèse b) entraîne que pour
$m\leq n$, $A_{im}$ est isomorphe à
$A_{in}/\mathfrak{K}^{m+1}A_{in}$. Par suite, le préschéma formel induit sur
$U_i$ par $\mathfrak{X}$ est isomorphe à $\operatorname{Spf}(A_i)$, où
$A_i=\varprojlim_n A_{in}$ (\hyperref[I.10.6.4-fr]{10.6.4})~; $A_i$ est un
anneau $\mathfrak{K}A_i$-adique, et $A_i/\mathfrak{K}A_i$, isomorphe à
$A_{i0}$, est une algèbre de type fini sur $B/\mathfrak{K}$. On en conclut
(\hyperref[0.7.5.5-fr]{0, 7.5.5}) que $A_i$ est topologiquement isomorphe à
un quotient d'une algèbre de séries formelles restreintes sur $B$ (par un
idéal nécessairement fermé, puisqu'une telle algèbre est noethérienne
(\hyperref[0.7.5.4-fr]{0, 7.5.4})).

Pour démontrer que c) entraîne a), on peut se limiter au cas où
$\mathfrak{X}=\operatorname{Spf}(A)$ est aussi affine, $A$ étant un anneau
adique noethérien, isomorphe au quotient d'une algèbre de séries formelles
restreintes sur $B$ par un idéal fermé. Alors
(\hyperref[0.7.5.5-fr]{0, 7.5.5}), $A/\mathfrak{K}A$ est une algèbre de type
fini sur $B/\mathfrak{K}$, et $\mathfrak{K}A=\mathfrak{J}$ est un idéal de
définition de $A$, donc, en vertu de
(\hyperref[I.10.10.9-fr]{10.10.9}), les conditions de a) sont satisfaites.

On notera que si les conditions de la prop.
(\hyperref[I.10.13.1-fr]{10.13.1}) sont remplies, la propriété a) est
valable pour \emph{tout Idéal de définition $\mathscr{K}$} de
$\mathfrak{Y}$ (en vertu de c)), et par suite, dans la propriété b),
\emph{tous les $f_n$ sont des morphismes de type fini}.

\begin{corollary}[10.13.2]
\label{I.10.13.2-fr}
Si les conditions de (\hyperref[I.10.13.1-fr]{10.13.1}) sont vérifiées,
tout ouvert formel affine noethérien $V$ de $\mathfrak{Y}$ possède la
propriété (\textbf{Q}) et si $\mathfrak{Y}$ est noethérien, il en est de
même de $\mathfrak{X}$.
\end{corollary}

Cela résulte aussitôt de (\hyperref[I.10.13.1-fr]{10.13.1}) et de
(\hyperref[I.6.3.2-fr]{6.3.2}).

\begin{definition}[10.13.3]
\label{I.10.13.3-fr}
Lorsque les propriétés équivalentes a), b), c) de
(\hyperref[I.10.13.1-fr]{10.13.1}) sont vérifiées, on dit que le morphisme
$f$ est \emph{de type fini}, ou que $\mathfrak{X}$ est un
\emph{$\mathfrak{Y}$-préschéma formel de type fini}, ou un
\emph{préschéma formel de type fini au-dessus de $\mathfrak{Y}$}.
\end{definition}

\begin{corollary}[10.13.4]
\label{I.10.13.4-fr}
Soient $\mathfrak{X}=\operatorname{Spf}(A)$,
$\mathfrak{Y}=\operatorname{Spf}(B)$ deux schémas formels affines
noethériens~; pour que $\mathfrak{X}$ soit de type fini sur
$\mathfrak{Y}$, il faut et il suffit que l'anneau adique noethérien $A$ soit
isomorphe au quotient d'une algèbre de séries formelles restreintes sur $B$
par un idéal fermé.
\end{corollary}

En effet, avec les notations de (\hyperref[I.10.13.1-fr]{10.13.1}), si
$\mathfrak{X}$ est de type fini sur $\mathfrak{Y}$,
$A/\mathfrak{K}A$ est alors une $(B/\mathfrak{K})$-algèbre de type fini en
vertu de (\hyperref[I.6.3.3-fr]{6.3.3}) et $\mathfrak{K}A$ est un idéal de
définition de $A$ (\hyperref[I.10.10.9-fr]{10.10.9}). On conclut donc par
(\hyperref[0.7.5.5-fr]{0, 7.5.5}).

\begin{proposition}[10.13.5]
\label{I.10.13.5-fr}
\medskip\noindent
\begin{enumerate}
  \item[\rm{(i)}] Le composé de deux morphismes de préschémas formels qui
    sont de type fini est de type fini.

